\documentclass{NoteThesis/note-thesis}

\course{Polymer Processing}

\coursetitle{QXU6018 Polymer Processing}
\semester{2025B}
\teacher{Prof. Xuetao Shi, Prof. Nan Tian, Prof. Zhenhua Wang}

\begin{document}

\maketitle

\tableofcontents

\newpage

\section{Introduction to Module}

\subsection*{Overview}

In spring semester, theory lecture.

In autumn semester, 6 lab experience.

\subsection*{Assessment}

\begin{table}[h]
    \centering
    \begin{tabular}{cc}
        \hline
        Activity & Score\\
        \hline
        Regular Homework Assignments & 10\%\\
        Group Presentation & 10\%\\
        Individual Coursework & 40\%\\
        Exam in Week10 & 20\%\\ 
        \hline
    \end{tabular}
\end{table}

\section{Mechanical Behavior of Polymers}

\subsection{Viscoelastic Properties of Polymers}

At a specific temperature and molecular weight, may behave as a liquid or a solid depend on time.

\subsubsection*{Comparison of Elastic vs. Viscous Behavior}

\begin{table}[h]
    \centering
    \label{tab:elastic_viscous}
    \begin{tabular}{lcc}
        \hline
        \textbf{Property} & \textbf{Elastic (Hookean)} & \textbf{Viscous (Newtonian)} \\
        \hline
        Constitutive Equation & $\sigma = E\epsilon$ & $\sigma = \mu\dot{\gamma}$ \\
        Response to Load & Instantaneous & Time-dependent \\
        Energy & Stored & Dissipated \\
        Recovery & Complete & None (permanent) \\
        Rate Dependence & No & Yes \\
        \hline
    \end{tabular}
\end{table}

\subsubsection*{Stress Relaxation}

The stress response to a unit constant strain:

\begin{equation*}
    E_e(t) = \sigma(t) / \epsilon(0)
\end{equation*}

\subsection{Time-Temperature Superposition(TTS)}

\subsubsection*{WLF Equation}

\begin{equation*}
    \log a_T = \frac{-C_1(T-T_{ref})}{C_2+T-T_{ref}}
\end{equation*}

$C_1$ and $C_2$ is material dependent constant.

If $T_g$ is the ref temperature, $C_1=17.44$ and $C_2=51.6$.

\subsubsection*{Assumptions}

\begin{enumerate}
    \item The material should not undergo any chemical or physical changes as result of the temperature change.
    \item No heterogeneity in sample, for example, blend.
    \item Application in linear Viscoelastic regime only.
\end{enumerate}

\subsection{Dynamic Mechanical Test}

\subsubsection*{General form}

1

\subsection{Fiber Reinforced Polymer}

1

\section{Rheology}

\subsection{What is Rheology}

Rheology is the study of flow and deformation of all forms of matter.

\[
N_1 = \tau_{xx} - \tau_{yy} = \psi_1 (T,\dot{\gamma}) \dot{\gamma_{xy}}^2
\]

\[
N_2 = \tau_{yy} - \tau_{zz} = \psi_2 (T,\dot{\gamma}) \dot{\gamma_{xy}}^2
\] 

\subsubsection*{Deborah Number}

\[
De = \frac{\lambda}{t_{process}}
\]

\[\cdots \]

\subsection{}

\end{document}
